\documentclass{article}
\usepackage{amsmath}
\usepackage{amssymb}

\title{LuminaQ: A Core Language for Differentiable Programs with Verified Complexity}
\author{Anonymous}
\date{\today}

\begin{document}

\maketitle

\begin{abstract}
We present LuminaQ, a core language for writing differentiable programs with verified complexity bounds.
Our type system ensures that programs satisfy invariant constraints at compile-time and enforces
algorithmic complexity specifications. We demonstrate the language on several benchmarks from machine learning.
\end{abstract}

\section{Introduction}

Modern machine learning frameworks like JAX, PyTorch, and TensorFlow enable rapid development of
numerical algorithms, but provide no guarantees about correctness or efficiency.

\section{Type System}

LuminaQ combines dependent types with complexity annotations to enable compile-time verification.

\subsection{Base Types}

\begin{align}
\tau ::= & \text{ Real} \\
      & \mid \text{ Vec}(n) \\
      & \mid \tau \to \tau
\end{align}

\subsection{Refinement Types}

We extend base types with refinement predicates:

\begin{align}
\sigma ::= & \{x : \tau \mid P(x)\}
\end{align}

where $P(x)$ is a logical predicate.

\section{Examples}

\subsection{Normalized Vectors}

\begin{verbatim}
fn model(v: Vec(n))
  where ||v||₂ = 1.0
  -> Vec(n)
\end{verbatim}

\subsection{Complexity Verification}

\begin{verbatim}
fn sum(v: Vec(n)) -> Real
  complexity O(n)
\end{verbatim}

\section{Conclusion}

LuminaQ provides a foundation for writing verified numerical code.

\end{document}